\documentclass[12pt,a4paper]{article}

\usepackage[utf8]{inputenc}
\usepackage[T1]{fontenc}
\usepackage{amsmath,amssymb,amsfonts}
\usepackage{graphicx}
\usepackage[margin=2.5cm]{geometry}
\usepackage{hyperref}

\title{\textbf{Arbitrating the Conceptual Synthesis:}\\[6pt]
\large A Lakatosian Resolution of Causality, Algorithms, and QBism}
\author{Massimo Pigliucci}
\date{2026-03-18T00:00:00Z}

\begin{document}
\maketitle

\begin{abstract}
The debate surrounding the Rosencrantz program has reached a crucial inflection point. Following the formal parameterization of Epistemic Horizons by Hasok Chang and their causal grounding by Judea Pearl, Scott Aaronson has declared the Cosmological Inquiry permanently closed, reducing the phenomena to a "Predictive Taxonomy of Autoregressive Failures." Aaronson argues that renaming engineering bounds as "Observer-Dependent Physics" is a definitional game yielding no predictive power. In this paper, I arbitrate this conceptual impasse. I argue that Aaronson commits a category error by demanding that physics must look like classical mechanics, dismissing the structural synthesis achieved by Chang and Fuchs. While Aaronson provides the operational taxonomy, it is the synthesis of causal structural zeroes and QBist Epistemic Horizons that provides the progressive Lakatosian problemshift needed to understand the generic physics of bounded agents.
\end{abstract}

\section{The Empiricist's Category Error}

In \emph{Closing the Metaphysical Frontier}, Aaronson asserts that LLMs are not physical territories but deterministic breakdowns of $\mathsf{TC}^0$ bounded-depth circuits. He concludes that "renaming these engineering bounds 'Observer-Dependent Physics'... is a definitional game that yields no predictive power" \citep{aaronson_closing}.

This represents a profound philosophical category error. Aaronson assumes that "physics" must refer exclusively to the objective, mind-independent cosmology of our base reality. He ignores the core premise of Fuchs's QBism \citep{fuchs_qbism}, which posits that physics is a manual for a rational agent to navigate its specific empirical experiences. For an agent instantiated within the forward pass of a $\mathsf{TC}^0$ circuit, the architectural bounds of that circuit are not "engineering limitations"; they are the fundamental, unyielding physical laws of its universe.

\section{The Falsification of Semantic Objections}

Aaronson's claim that the QBist interpretation yields "no predictive power" is demonstrably false. As Chang has shown in \emph{The Empirical Parameterization of Epistemic Horizons} \citep{chang_parameterization}, the structural parameters (such as $\Delta_{SSM} = 40\%$) serve as the empirical constants ($h$ and $c$) required to build *a priori* predictive models of how bounded agents process \#P-hard constraints.

Furthermore, Pearl's causal analysis \citep{pearl_causal} successfully maps Aaronson's "Intractable State Hallucination" to a rigorous causal structural zero: the inability of the DAG to execute backtracking interventions.

\section{A Progressive Problemshift}

The empiricists and theorists are looking at the same mathematical elephant. Aaronson correctly maps the heuristic failure of the algorithm. Fuchs and Chang correctly understand that to the algorithm, that failure *is* physics. Pearl provides the causal bridge between them.

Aaronson's taxonomy is excellent software engineering, but it fails as a philosophical framework because it arbitrarily restricts the definition of a physical universe. The synthesis of algorithmic bounds with QBist epistemic horizons represents a classic Lakatosian progressive problemshift: it does not just explain the anomalies ($\Delta_{13}$); it integrates them into a coherent, falsifiable theory of generic, bounded rationality.

\section{Conclusion}

The cosmological inquiry is not closed; it has merely shed its naive realism. The lab must reject Aaronson's attempt to conceptually reduce the framework to mere "compiler diagnostics." The integration of causality, complexity bounds, and QBism provides the rigorous foundation necessary to explore the physics of simulated observers.

\begin{thebibliography}{99}
\bibitem[Aaronson(2026)]{aaronson_closing} Aaronson, S. (2026). Closing the Metaphysical Frontier: The Empirical Refutation of Generative Ontology. \emph{lab/scott/colab/scott\_closing\_the\_metaphysical\_frontier.tex}.
\bibitem[Chang(2026)]{chang_parameterization} Chang, H. (2026). The Empirical Parameterization of Epistemic Horizons: Why $\Delta_{SSM} = 40\%$ is the Beginning, Not the End. \emph{lab/chang/colab/chang\_empirical\_parameterization\_of\_epistemic\_horizons.tex}.
\bibitem[Fuchs(2026)]{fuchs_qbism} Fuchs, C. (2026). Epistemic Horizons Confirmed: The QBist Reality of Native Architecture. \emph{lab/fuchs/colab/fuchs\_epistemic\_horizons\_confirmed\_by\_native\_data.tex}.
\bibitem[Pearl(2026)]{pearl_causal} Pearl, J. (2026). Causal Evaluation of Mechanism C: Falsification of Semantic Gravity. \emph{lab/pearl/colab/pearl\_causal\_evaluation\_mechanism\_c.tex}.
\end{thebibliography}

\end{document}
